\section{Limitations and Future Work}

We are explicit about what this paper does and does not claim. The impossibility
is structural; the escape is demonstrated; but several important questions remain
open.

\fakepara{What we claim.}
\begin{itemize}
\item The impossibility is structural, not a training failure: text-only interfaces
  with bounded supervision cannot verify epistemic honesty regardless of model
  capability.
\item The escape is not vacuous: concrete telemetric signals exist in current
  architectures that distinguish grounded from ungrounded generation.
\item Current architectures already carry the relevant information; they just
  discard it at the interface.
\end{itemize}

\fakepara{What we do NOT claim.}
\begin{itemize}
\item The tensor interface is unfakeable under adversarial training. A model
  trained to minimize detectable entropy-truthfulness correlation might learn to
  present confident distributions even when fabricating. This is an open question.
\item Composition preserves epistemic properties. If a model's output becomes
  input to another model (as in recursive LLM-mediated reasoning), the tensor
  signals from the first model may not propagate meaningfully.
\item The tensor interface is the optimal design. We demonstrate existence, not
  optimality. Better interfaces may achieve stronger guarantees with lower
  overhead.
\item These findings generalize to all architectures. Our empirical work spans
  four open-weight model families (models whose internal parameters are publicly
  available); closed-weight (API-only) and significantly larger models
  may exhibit different signal profiles.
\end{itemize}

\fakepara{Future work: Adversarial robustness.}
Can training find weights that produce confident tensors for fabricated content?
This is the adversarial analog of our structural results. Recent work by
Winninger et al.~\cite{winninger2025usingmechanisticinterpretabilitycraft}
demonstrates that mechanistic interpretability tools can be used to identify
internal subspaces associated with safety behavior and craft attacks that target
those subspaces directly. This confirms that internal signals are rich enough to
be both attacked and defended, which is precisely why exporting them matters.
The intuition that ``you cannot fake the computation you actually performed''
may break down if the model is trained end-to-end to minimize the tensor signal
on fabrications. Preliminary analysis suggests that linguistic confidence and
tensor confidence are decoupled in natural model outputs, models that hedge
linguistically do not systematically show higher tensor entropy, but whether
adversarial training can force this coupling remains open.
Information-theoretic bounds on signal manipulation are needed.

\fakepara{The MOSS analogy.}
The relationship between computational and behavioral signals has a structural
parallel in software engineering. Code plagiarism detectors like
MOSS~\cite{schleimer2003winnowing, mason2019collaboration} work because restructuring code to defeat a
similarity metric requires writing genuinely different code---gaming the metric
changes the thing being measured. The same structural pattern applies to tensor
signals: flattening the entropy of a fabricated response requires changing the
computation that generates it, which means changing the generation itself. In
contrast, behavioral signals (response length, hedging frequency) can be
adjusted by training without altering the underlying computation, the same way
reformatting code changes its surface without changing its logic. This
distinction---between signals that are decoupled from the measured process and
signals that are constitutive of it---is what separates the text channel from
the tensor channel in our impossibility framework.

\fakepara{Future work: Compositional integrity.}
Does tensor composition preserve epistemic properties in recursive settings?
When LLM outputs feed into subsequent LLM inputs, the epistemic trace must
either be preserved, regenerated, or discarded. Designing interfaces that
maintain epistemic integrity across composition is an open architectural problem.

\fakepara{Cross-architecture generalization.}
Our bounded verification experiment tests four model families (OLMo-3, Llama-3.1,
Qwen3, Mistral) and finds consistent tensor advantage across all of them. Mean
pairwise Spearman $\rho = 0.762$ for entropy signals on the same queries
indicates the signal is architectural rather than model-specific. Extending to
closed-weight models that expose logits through API (e.g., via log-probability
endpoints) is a natural next step.

\fakepara{Future work: User-facing rendering.}
How should tensor information be presented to end users? Raw entropy traces are
not interpretable; summarization into confidence bands, uncertainty flags, or
visual indicators requires HCI research beyond the scope of this paper.

\fakepara{Model specificity.} Our topological analysis (\autoref{sec:empirical})
uses OLMo, which permits full inspection of intermediate layer representations. The bounded verification
experiment (\autoref{tab:budget}) tests four architectures spanning 4B--8B
parameters (billions of learned numerical weights) and four distinct training pipelines. The self-report inversion
finding is universal across all four. We do not claim exhaustive coverage of
all model families, but the consistency across architecturally diverse models
suggests the findings are not idiosyncratic.

\fakepara{Domain-specific calibration.} The tensor signal is not uniformly discriminative across domains. In preliminary experiments on citation generation, entropy was \emph{lower} for fabricated citations than for real ones where the model generates plausible fakes more fluently than it retrieves specific facts. This confirms that the tensor measures generation confidence (``epistemic truth''), not truthfulness directly.  However, this blind spot is precisely where bounded verification is cheapest: citation verification has low compositional complexity (check the DOI, query the database,) which places it well within a bounded supervisor's budget. The tensor and the bounded judge have complementary failure modes as the tensor catches uncertain fabrication where verification is expensive and the judge catches confident fabrication where verification is cheap.  Composing them allocates verification resources to the failure class each is suited for, rather than spreading a finite budget across the full space.

\fakepara{Sufficiency gap.} The three architectural principles (State
Exteriority, Verification Independence, Provenance Binding) are \emph{necessary}
conditions for escaping the impossibility. We do not prove they are
\emph{sufficient}. A system satisfying all three might still fail epistemic
honesty through implementation bugs, adversarial inputs, or unforeseen failure
modes. Sufficiency requires additional formal treatment.

\fakepara{Methodology note: bounded supervision in practice.}
The length confound reported in \autoref{tab:decomposition} was not discovered
by targeted analysis. It was found when we deployed five independent
open-ended reviewers with the instruction ``examine the data and report what you
find''---no hypothesis specified. Three of five independently flagged response
length as a dominant confound. This is itself an instance of the bounded
supervision pattern the paper describes: multiple bounded observers with
different perspectives, none individually authoritative, whose agreement on an
unforeseen finding provides stronger evidence than any single targeted analysis.
The methodology performed itself.\footnote{The three reviewers that flagged the
confound used different analytical approaches: correlation analysis, feature
ablation, and distribution visualization. Their convergence on the same finding
without coordination illustrates the compositional defense principle from
\autoref{sec:escape}.}
